% Panorama de MCA_WINDOW EXPANDIDO, a nivel de operación, en el estilo de
% event_window_capture del OSC: el disparo y el asignador son compuertas; las
% marcas, contadores con su +1; la muestra, una resta y un máximo; la
% contabilidad, compuertas y contadores con su habilitación. Las ranuras y el
% banco de trapecios quedan como submódulos.
% Las redes con nombre (disparo, veto, puede, esperan, abre...) se conectan por
% nombre: el trazo corto que entra a un pin lleva el nombre de la red. Las que
% vienen de las ranuras (libre, abierta, cerrada, cierra...) son vectores de un
% bit por ranura.
% Fuente: modulos/mca/rtl/mca_window.sv, instanciado por mca_pulse_feature como
% en rtl/mine/tops/mca_red_pitaya_top.sv (3 ranuras, 3 formas).
% Estilo: docs/figuras/estilo_tikz.tex; los colores por función, los mismos de
% panorama_mca_top.tex. El borde grueso es el límite del módulo.
% El epígrafe no va en la figura: está en panorama_mca_window_expandido.txt.
%
% Compilar:   pdflatex panorama_mca_window_expandido.tex
%             pdftocairo -svg panorama_mca_window_expandido.pdf panorama_mca_window_expandido.svg
%
\documentclass[border=0pt,tikz]{standalone}
\usepackage[utf8]{inputenc}
\usepackage[T1]{fontenc}
\usepackage{lmodern}
\usepackage[none]{hyphenat}
% Estilo común de las figuras TikZ (panoramas y diagramas de bloques).
%
% Cada figura lo incluye con % Estilo común de las figuras TikZ (panoramas y diagramas de bloques).
%
% Cada figura lo incluye con % Estilo común de las figuras TikZ (panoramas y diagramas de bloques).
%
% Cada figura lo incluye con \input{../../figuras/estilo_tikz.tex} (las figuras
% viven en docs/<módulo>/figuras/). Un solo lugar para la paleta, los estilos de
% bloque y el formato de los parámetros: si cambia algo acá, cambia en todas.
%
% La paleta es la misma de docs/figuras/estilo_svg.py y de PALETA en
% modulos/mk/datasheet.py; si se toca un color, se toca en los tres.
%
% Requiere: \usepackage{tikz} (lo da la clase standalone con la opción tikz).

\usetikzlibrary{arrows.meta,positioning,calc,fit,backgrounds,shapes.multipart,
                decorations.markings,shadows}

% ---------------------------------------------------------------- colores ---
\definecolor{azul}{RGB}{108,172,228}      % acondicionamiento
\definecolor{verde}{RGB}{110,214,142}     % captura / histogramas en BRAM
\definecolor{amarillo}{RGB}{252,232,160}  % trigger
\definecolor{naranja}{RGB}{252,206,160}   % camino a DDR
\definecolor{grisr}{RGB}{226,226,226}     % compartido / infraestructura
\definecolor{rojo}{RGB}{240,128,128}      % bus del sistema

% ---------------------------------------------------------------- estilos ---
\tikzset{
  >={Stealth[length=2.4mm,width=1.8mm]},
  % fondo blanco siempre: el SVG se ve igual en modo claro y oscuro
  fondo/.style   = {show background rectangle, inner frame sep=8pt,
                    background rectangle/.style={fill=white, draw=none}},
  bloque/.style  = {draw=black, line width=0.6pt, rounded corners=2pt,
                    align=center, text width=2.2cm, minimum height=1.15cm,
                    inner sep=2pt, font=\scriptsize, fill=#1},
  datos/.style   = {bloque=azul!35},
  trig/.style    = {bloque=amarillo},
  captura/.style = {bloque=verde!55},
  ddr/.style     = {bloque=naranja!70},
  comp/.style    = {bloque=grisr},
  bus/.style     = {bloque=rojo!45},
  % lógica escrita dentro del módulo, sin submódulo propio
  etapa/.style   = {bloque=#1, line width=0.5pt},
  etapa/.default = verde!20,
  % lo que el bitstream no trae (atado a cero, stub, salida muda)
  mudo/.style    = {bloque=white, draw=black!45, dashed, text=black!55},
  %
  % BLOQUE CON PARÁMETROS, en el formato del datasheet: el encabezado lleva el
  % nombre y los parámetros con que lo instancia el padre (NOMBRE=valor, en
  % monoespaciada, uno por línea); el cuerpo, más claro, lo que hace.
  %   \node[datos, param={azul!35}{azul!12}] {\blk{nombre}{} \prm{A=1\\B=2}
  %                                            \nodepart{two} función};
  param/.style 2 args = {rectangle split, rectangle split parts=2,
                         rectangle split part fill={#1, #2}},
  %
  flujo/.style   = {->, line width=0.8pt},
  ctrl/.style    = {->, line width=0.6pt, black!65},
  % la línea que pasa por encima en un cruce: un halo blanco la separa
  cruce/.style   = {preaction={draw=white, line width=3.4pt, -}},
  % marca de ancho de bus: barra oblicua a mitad de la flecha, con el número
  % encima. \draw[flujo, ancho=14] ...; para otra posición, ancho en=0.3.
  ancho/.style   = {decoration={markings, mark=at position \anchopos with {
                      \draw[line width=0.6pt, -] (-0.1,-0.14) -- (0.1,0.14);
                      \node[font=\scriptsize, above=1.5pt] {#1};}},
                    postaction={decorate}},
  % la misma marca en una flecha vertical: el número va a la derecha
  anchov/.style  = {decoration={markings, mark=at position \anchopos with {
                      \draw[line width=0.6pt, -] (-0.14,-0.1) -- (0.14,0.1);
                      \node[font=\scriptsize, right=3pt] {#1};}},
                    postaction={decorate}},
  ancho en/.code = {\def\anchopos{#1}},
  sen/.style     = {font=\scriptsize\ttfamily, inner sep=1.5pt},
  nodo/.style    = {font=\tiny\ttfamily, inner sep=1pt, black!75},
  pin/.style     = {font=\scriptsize, align=center, inner sep=1pt},
  grupo/.style   = {draw=black!55, dashed, line width=0.7pt, rounded corners=4pt},
  canal/.style   = {grupo},
  contenedor/.style = {draw=black!80, line width=1.0pt, rounded corners=6pt},
  % marca numerada: cada número es una fila de la tabla del documento
  marca/.style   = {circle, fill=black, text=white, font=\scriptsize\bfseries,
                    inner sep=0pt, minimum size=4.3mm},
  %
  % DECISIONES (docs/mca/vista_general_mca.md). Cada opción es un bloque, todas
  % entran a un mux (\muxdec, abajo) y el registro que elige entra por abajo.
  % la opción que trae el reset (el comportamiento histórico): borde doble
  defecto/.style = {double=white, double distance=0.8pt, line width=0.45pt},
  % rótulo de la decisión, el mismo Dn del panorama y del documento
  decision/.style = {circle, fill=black, text=white, font=\small\bfseries,
                     inner sep=0pt, minimum size=7mm},
  %
  % LÓGICA (docs/mca/figuras/contabilidad/ y los detalles de las decisiones)
  % registro: caja gris con el triángulo de reloj en el borde de abajo
  registro/.style = {draw=black, line width=0.6pt, fill=grisr, align=center,
                     font=\scriptsize, inner sep=2pt, minimum width=#1,
                     minimum height=0.85cm,
                     path picture={\draw[line width=0.5pt]
                       ([xshift=-1.1mm]path picture bounding box.south) --
                       ([yshift=1.6mm]path picture bounding box.south) --
                       ([xshift=1.1mm]path picture bounding box.south);}},
  registro/.default = 0.8cm,
  % condición de una cadena if / else-if: "sí" sale a la derecha, "no" abajo
  pregunta/.style = {draw=black, line width=0.5pt, rounded corners=7pt,
                     fill=white, font=\scriptsize\ttfamily, align=center,
                     inner sep=2.5pt, minimum height=0.62cm, text width=#1},
  pregunta/.default = 1.9cm,
  si/.style      = {font=\tiny, inner sep=1pt, above, pos=0.12},
  no/.style      = {font=\tiny, inner sep=1pt, left, pos=0.45},
  %
  titulo/.style  = {anchor=north west, font=\small, inner sep=4pt, align=left},
  rotulo/.style  = {anchor=north west, font=\scriptsize\itshape, black!70,
                    inner sep=2pt, align=left},
}

% submódulo: nombre en monoespaciada; lógica propia: nombre en romana
\newcommand{\blk}[2]{{\ttfamily\bfseries #1}\\[1pt]#2}
\newcommand{\etq}[2]{{\bfseries #1}\\[1pt]#2}
\newcommand{\punto}[1]{\fill (#1) circle (1.3pt);}

% Parámetros en el formato del datasheet: NOMBRE=valor en monoespaciada, una
% entrada por celda, hasta cuatro columnas. Una columna: \prm{A=1\\B=2};
% dos: \prm{A=1 & B=2\\C=3}.
\newcommand{\prm}[1]{{\ttfamily\fontsize{5.6}{6.6}\selectfont
  \begin{tabular}{@{}l@{\hspace{7pt}}l@{\hspace{7pt}}l@{\hspace{7pt}}l@{}}#1\end{tabular}}}

\def\anchopos{0.5}

% Mux de una decisión: trapecio con el lado ancho (las entradas) a la izquierda.
% #1 = nombre, #2 = centro, #3 = alto en cm. Deja tres coordenadas:
%   (#1-in)  centro del lado de entradas: las flechas llegan a (#1-in |- opcion.east)
%   (#1-out) la salida
%   (#1-sel) mitad del lado de abajo: por ahí entra el registro que elige
% El rótulo de adentro es opcional: \muxdec{m}{(4,0)}{3.6} y después
% \node[...] at (m) {...}.
\newcommand{\muxdec}[3]{%
  \path #2 coordinate (#1);
  \draw[line width=0.6pt, fill=grisr]
    ($(#1)+(-0.45,{#3/2})$) -- ($(#1)+(0.45,{#3/2-0.35})$) --
    ($(#1)+(0.45,{-#3/2+0.35})$) -- ($(#1)+(-0.45,{-#3/2})$) -- cycle;
  \coordinate (#1-in)  at ($(#1)+(-0.45,0)$);
  \coordinate (#1-out) at ($(#1)+(0.45,0)$);
  \coordinate (#1-sel) at ($(#1)+(0,{-#3/2+0.175})$);}

% El mismo mux rotado 90 grados, para un flujo de arriba hacia abajo: entradas
% por el lado ancho de arriba, salida abajo y el registro que elige por la
% izquierda. #1 = nombre, #2 = centro, #3 = ancho en cm. Deja (#1-in) centro del
% lado de entradas, (#1-out) la salida y (#1-sel) mitad del lado izquierdo.
\newcommand{\muxdecv}[3]{%
  \path #2 coordinate (#1);
  \draw[line width=0.6pt, fill=grisr]
    ($(#1)+({-#3/2},0.45)$) -- ($(#1)+({#3/2},0.45)$) --
    ($(#1)+({#3/2-0.35},-0.45)$) -- ($(#1)+({-#3/2+0.35},-0.45)$) -- cycle;
  \coordinate (#1-in)  at ($(#1)+(0,0.45)$);
  \coordinate (#1-out) at ($(#1)+(0,-0.45)$);
  \coordinate (#1-sel) at ($(#1)+({-#3/2+0.175},0)$);}

% Leyenda de colores + nota. #1 = posición, #2 = color/texto, #3 = nota.
\newcommand{\leyenda}[3]{%
  \begin{scope}[shift={(#1)}, every node/.style={font=\scriptsize}]
    \coordinate (t0) at (-0.5,0);
    \foreach \c/\t [count=\i] in {#2} {
      \pgfmathtruncatemacro\j{\i-1}
      \node[draw, fill=\c, minimum width=4mm, minimum height=3mm, inner sep=0pt,
            right=5mm of t\j.east, anchor=west] (l\i) {};
      \node[right=1.5mm of l\i] (t\i) {\t};
    }
    \node[anchor=north west, black!70, text width=19cm, align=left]
      at (-0.05,-0.3) {#3};
  \end{scope}}

% Lo que dicen los parámetros, para la nota de la leyenda
\newcommand{\notaparam}{\texttt{NOMBRE=valor}: parámetros con que el padre
  instancia el bloque (se omiten los que sólo pasan el del padre sin cambio).}
 (las figuras
% viven en docs/<módulo>/figuras/). Un solo lugar para la paleta, los estilos de
% bloque y el formato de los parámetros: si cambia algo acá, cambia en todas.
%
% La paleta es la misma de docs/figuras/estilo_svg.py y de PALETA en
% modulos/mk/datasheet.py; si se toca un color, se toca en los tres.
%
% Requiere: \usepackage{tikz} (lo da la clase standalone con la opción tikz).

\usetikzlibrary{arrows.meta,positioning,calc,fit,backgrounds,shapes.multipart,
                decorations.markings,shadows}

% ---------------------------------------------------------------- colores ---
\definecolor{azul}{RGB}{108,172,228}      % acondicionamiento
\definecolor{verde}{RGB}{110,214,142}     % captura / histogramas en BRAM
\definecolor{amarillo}{RGB}{252,232,160}  % trigger
\definecolor{naranja}{RGB}{252,206,160}   % camino a DDR
\definecolor{grisr}{RGB}{226,226,226}     % compartido / infraestructura
\definecolor{rojo}{RGB}{240,128,128}      % bus del sistema

% ---------------------------------------------------------------- estilos ---
\tikzset{
  >={Stealth[length=2.4mm,width=1.8mm]},
  % fondo blanco siempre: el SVG se ve igual en modo claro y oscuro
  fondo/.style   = {show background rectangle, inner frame sep=8pt,
                    background rectangle/.style={fill=white, draw=none}},
  bloque/.style  = {draw=black, line width=0.6pt, rounded corners=2pt,
                    align=center, text width=2.2cm, minimum height=1.15cm,
                    inner sep=2pt, font=\scriptsize, fill=#1},
  datos/.style   = {bloque=azul!35},
  trig/.style    = {bloque=amarillo},
  captura/.style = {bloque=verde!55},
  ddr/.style     = {bloque=naranja!70},
  comp/.style    = {bloque=grisr},
  bus/.style     = {bloque=rojo!45},
  % lógica escrita dentro del módulo, sin submódulo propio
  etapa/.style   = {bloque=#1, line width=0.5pt},
  etapa/.default = verde!20,
  % lo que el bitstream no trae (atado a cero, stub, salida muda)
  mudo/.style    = {bloque=white, draw=black!45, dashed, text=black!55},
  %
  % BLOQUE CON PARÁMETROS, en el formato del datasheet: el encabezado lleva el
  % nombre y los parámetros con que lo instancia el padre (NOMBRE=valor, en
  % monoespaciada, uno por línea); el cuerpo, más claro, lo que hace.
  %   \node[datos, param={azul!35}{azul!12}] {\blk{nombre}{} \prm{A=1\\B=2}
  %                                            \nodepart{two} función};
  param/.style 2 args = {rectangle split, rectangle split parts=2,
                         rectangle split part fill={#1, #2}},
  %
  flujo/.style   = {->, line width=0.8pt},
  ctrl/.style    = {->, line width=0.6pt, black!65},
  % la línea que pasa por encima en un cruce: un halo blanco la separa
  cruce/.style   = {preaction={draw=white, line width=3.4pt, -}},
  % marca de ancho de bus: barra oblicua a mitad de la flecha, con el número
  % encima. \draw[flujo, ancho=14] ...; para otra posición, ancho en=0.3.
  ancho/.style   = {decoration={markings, mark=at position \anchopos with {
                      \draw[line width=0.6pt, -] (-0.1,-0.14) -- (0.1,0.14);
                      \node[font=\scriptsize, above=1.5pt] {#1};}},
                    postaction={decorate}},
  % la misma marca en una flecha vertical: el número va a la derecha
  anchov/.style  = {decoration={markings, mark=at position \anchopos with {
                      \draw[line width=0.6pt, -] (-0.14,-0.1) -- (0.14,0.1);
                      \node[font=\scriptsize, right=3pt] {#1};}},
                    postaction={decorate}},
  ancho en/.code = {\def\anchopos{#1}},
  sen/.style     = {font=\scriptsize\ttfamily, inner sep=1.5pt},
  nodo/.style    = {font=\tiny\ttfamily, inner sep=1pt, black!75},
  pin/.style     = {font=\scriptsize, align=center, inner sep=1pt},
  grupo/.style   = {draw=black!55, dashed, line width=0.7pt, rounded corners=4pt},
  canal/.style   = {grupo},
  contenedor/.style = {draw=black!80, line width=1.0pt, rounded corners=6pt},
  % marca numerada: cada número es una fila de la tabla del documento
  marca/.style   = {circle, fill=black, text=white, font=\scriptsize\bfseries,
                    inner sep=0pt, minimum size=4.3mm},
  %
  % DECISIONES (docs/mca/vista_general_mca.md). Cada opción es un bloque, todas
  % entran a un mux (\muxdec, abajo) y el registro que elige entra por abajo.
  % la opción que trae el reset (el comportamiento histórico): borde doble
  defecto/.style = {double=white, double distance=0.8pt, line width=0.45pt},
  % rótulo de la decisión, el mismo Dn del panorama y del documento
  decision/.style = {circle, fill=black, text=white, font=\small\bfseries,
                     inner sep=0pt, minimum size=7mm},
  %
  % LÓGICA (docs/mca/figuras/contabilidad/ y los detalles de las decisiones)
  % registro: caja gris con el triángulo de reloj en el borde de abajo
  registro/.style = {draw=black, line width=0.6pt, fill=grisr, align=center,
                     font=\scriptsize, inner sep=2pt, minimum width=#1,
                     minimum height=0.85cm,
                     path picture={\draw[line width=0.5pt]
                       ([xshift=-1.1mm]path picture bounding box.south) --
                       ([yshift=1.6mm]path picture bounding box.south) --
                       ([xshift=1.1mm]path picture bounding box.south);}},
  registro/.default = 0.8cm,
  % condición de una cadena if / else-if: "sí" sale a la derecha, "no" abajo
  pregunta/.style = {draw=black, line width=0.5pt, rounded corners=7pt,
                     fill=white, font=\scriptsize\ttfamily, align=center,
                     inner sep=2.5pt, minimum height=0.62cm, text width=#1},
  pregunta/.default = 1.9cm,
  si/.style      = {font=\tiny, inner sep=1pt, above, pos=0.12},
  no/.style      = {font=\tiny, inner sep=1pt, left, pos=0.45},
  %
  titulo/.style  = {anchor=north west, font=\small, inner sep=4pt, align=left},
  rotulo/.style  = {anchor=north west, font=\scriptsize\itshape, black!70,
                    inner sep=2pt, align=left},
}

% submódulo: nombre en monoespaciada; lógica propia: nombre en romana
\newcommand{\blk}[2]{{\ttfamily\bfseries #1}\\[1pt]#2}
\newcommand{\etq}[2]{{\bfseries #1}\\[1pt]#2}
\newcommand{\punto}[1]{\fill (#1) circle (1.3pt);}

% Parámetros en el formato del datasheet: NOMBRE=valor en monoespaciada, una
% entrada por celda, hasta cuatro columnas. Una columna: \prm{A=1\\B=2};
% dos: \prm{A=1 & B=2\\C=3}.
\newcommand{\prm}[1]{{\ttfamily\fontsize{5.6}{6.6}\selectfont
  \begin{tabular}{@{}l@{\hspace{7pt}}l@{\hspace{7pt}}l@{\hspace{7pt}}l@{}}#1\end{tabular}}}

\def\anchopos{0.5}

% Mux de una decisión: trapecio con el lado ancho (las entradas) a la izquierda.
% #1 = nombre, #2 = centro, #3 = alto en cm. Deja tres coordenadas:
%   (#1-in)  centro del lado de entradas: las flechas llegan a (#1-in |- opcion.east)
%   (#1-out) la salida
%   (#1-sel) mitad del lado de abajo: por ahí entra el registro que elige
% El rótulo de adentro es opcional: \muxdec{m}{(4,0)}{3.6} y después
% \node[...] at (m) {...}.
\newcommand{\muxdec}[3]{%
  \path #2 coordinate (#1);
  \draw[line width=0.6pt, fill=grisr]
    ($(#1)+(-0.45,{#3/2})$) -- ($(#1)+(0.45,{#3/2-0.35})$) --
    ($(#1)+(0.45,{-#3/2+0.35})$) -- ($(#1)+(-0.45,{-#3/2})$) -- cycle;
  \coordinate (#1-in)  at ($(#1)+(-0.45,0)$);
  \coordinate (#1-out) at ($(#1)+(0.45,0)$);
  \coordinate (#1-sel) at ($(#1)+(0,{-#3/2+0.175})$);}

% El mismo mux rotado 90 grados, para un flujo de arriba hacia abajo: entradas
% por el lado ancho de arriba, salida abajo y el registro que elige por la
% izquierda. #1 = nombre, #2 = centro, #3 = ancho en cm. Deja (#1-in) centro del
% lado de entradas, (#1-out) la salida y (#1-sel) mitad del lado izquierdo.
\newcommand{\muxdecv}[3]{%
  \path #2 coordinate (#1);
  \draw[line width=0.6pt, fill=grisr]
    ($(#1)+({-#3/2},0.45)$) -- ($(#1)+({#3/2},0.45)$) --
    ($(#1)+({#3/2-0.35},-0.45)$) -- ($(#1)+({-#3/2+0.35},-0.45)$) -- cycle;
  \coordinate (#1-in)  at ($(#1)+(0,0.45)$);
  \coordinate (#1-out) at ($(#1)+(0,-0.45)$);
  \coordinate (#1-sel) at ($(#1)+({-#3/2+0.175},0)$);}

% Leyenda de colores + nota. #1 = posición, #2 = color/texto, #3 = nota.
\newcommand{\leyenda}[3]{%
  \begin{scope}[shift={(#1)}, every node/.style={font=\scriptsize}]
    \coordinate (t0) at (-0.5,0);
    \foreach \c/\t [count=\i] in {#2} {
      \pgfmathtruncatemacro\j{\i-1}
      \node[draw, fill=\c, minimum width=4mm, minimum height=3mm, inner sep=0pt,
            right=5mm of t\j.east, anchor=west] (l\i) {};
      \node[right=1.5mm of l\i] (t\i) {\t};
    }
    \node[anchor=north west, black!70, text width=19cm, align=left]
      at (-0.05,-0.3) {#3};
  \end{scope}}

% Lo que dicen los parámetros, para la nota de la leyenda
\newcommand{\notaparam}{\texttt{NOMBRE=valor}: parámetros con que el padre
  instancia el bloque (se omiten los que sólo pasan el del padre sin cambio).}
 (las figuras
% viven en docs/<módulo>/figuras/). Un solo lugar para la paleta, los estilos de
% bloque y el formato de los parámetros: si cambia algo acá, cambia en todas.
%
% La paleta es la misma de docs/figuras/estilo_svg.py y de PALETA en
% modulos/mk/datasheet.py; si se toca un color, se toca en los tres.
%
% Requiere: \usepackage{tikz} (lo da la clase standalone con la opción tikz).

\usetikzlibrary{arrows.meta,positioning,calc,fit,backgrounds,shapes.multipart,
                decorations.markings,shadows}

% ---------------------------------------------------------------- colores ---
\definecolor{azul}{RGB}{108,172,228}      % acondicionamiento
\definecolor{verde}{RGB}{110,214,142}     % captura / histogramas en BRAM
\definecolor{amarillo}{RGB}{252,232,160}  % trigger
\definecolor{naranja}{RGB}{252,206,160}   % camino a DDR
\definecolor{grisr}{RGB}{226,226,226}     % compartido / infraestructura
\definecolor{rojo}{RGB}{240,128,128}      % bus del sistema

% ---------------------------------------------------------------- estilos ---
\tikzset{
  >={Stealth[length=2.4mm,width=1.8mm]},
  % fondo blanco siempre: el SVG se ve igual en modo claro y oscuro
  fondo/.style   = {show background rectangle, inner frame sep=8pt,
                    background rectangle/.style={fill=white, draw=none}},
  bloque/.style  = {draw=black, line width=0.6pt, rounded corners=2pt,
                    align=center, text width=2.2cm, minimum height=1.15cm,
                    inner sep=2pt, font=\scriptsize, fill=#1},
  datos/.style   = {bloque=azul!35},
  trig/.style    = {bloque=amarillo},
  captura/.style = {bloque=verde!55},
  ddr/.style     = {bloque=naranja!70},
  comp/.style    = {bloque=grisr},
  bus/.style     = {bloque=rojo!45},
  % lógica escrita dentro del módulo, sin submódulo propio
  etapa/.style   = {bloque=#1, line width=0.5pt},
  etapa/.default = verde!20,
  % lo que el bitstream no trae (atado a cero, stub, salida muda)
  mudo/.style    = {bloque=white, draw=black!45, dashed, text=black!55},
  %
  % BLOQUE CON PARÁMETROS, en el formato del datasheet: el encabezado lleva el
  % nombre y los parámetros con que lo instancia el padre (NOMBRE=valor, en
  % monoespaciada, uno por línea); el cuerpo, más claro, lo que hace.
  %   \node[datos, param={azul!35}{azul!12}] {\blk{nombre}{} \prm{A=1\\B=2}
  %                                            \nodepart{two} función};
  param/.style 2 args = {rectangle split, rectangle split parts=2,
                         rectangle split part fill={#1, #2}},
  %
  flujo/.style   = {->, line width=0.8pt},
  ctrl/.style    = {->, line width=0.6pt, black!65},
  % la línea que pasa por encima en un cruce: un halo blanco la separa
  cruce/.style   = {preaction={draw=white, line width=3.4pt, -}},
  % marca de ancho de bus: barra oblicua a mitad de la flecha, con el número
  % encima. \draw[flujo, ancho=14] ...; para otra posición, ancho en=0.3.
  ancho/.style   = {decoration={markings, mark=at position \anchopos with {
                      \draw[line width=0.6pt, -] (-0.1,-0.14) -- (0.1,0.14);
                      \node[font=\scriptsize, above=1.5pt] {#1};}},
                    postaction={decorate}},
  % la misma marca en una flecha vertical: el número va a la derecha
  anchov/.style  = {decoration={markings, mark=at position \anchopos with {
                      \draw[line width=0.6pt, -] (-0.14,-0.1) -- (0.14,0.1);
                      \node[font=\scriptsize, right=3pt] {#1};}},
                    postaction={decorate}},
  ancho en/.code = {\def\anchopos{#1}},
  sen/.style     = {font=\scriptsize\ttfamily, inner sep=1.5pt},
  nodo/.style    = {font=\tiny\ttfamily, inner sep=1pt, black!75},
  pin/.style     = {font=\scriptsize, align=center, inner sep=1pt},
  grupo/.style   = {draw=black!55, dashed, line width=0.7pt, rounded corners=4pt},
  canal/.style   = {grupo},
  contenedor/.style = {draw=black!80, line width=1.0pt, rounded corners=6pt},
  % marca numerada: cada número es una fila de la tabla del documento
  marca/.style   = {circle, fill=black, text=white, font=\scriptsize\bfseries,
                    inner sep=0pt, minimum size=4.3mm},
  %
  % DECISIONES (docs/mca/vista_general_mca.md). Cada opción es un bloque, todas
  % entran a un mux (\muxdec, abajo) y el registro que elige entra por abajo.
  % la opción que trae el reset (el comportamiento histórico): borde doble
  defecto/.style = {double=white, double distance=0.8pt, line width=0.45pt},
  % rótulo de la decisión, el mismo Dn del panorama y del documento
  decision/.style = {circle, fill=black, text=white, font=\small\bfseries,
                     inner sep=0pt, minimum size=7mm},
  %
  % LÓGICA (docs/mca/figuras/contabilidad/ y los detalles de las decisiones)
  % registro: caja gris con el triángulo de reloj en el borde de abajo
  registro/.style = {draw=black, line width=0.6pt, fill=grisr, align=center,
                     font=\scriptsize, inner sep=2pt, minimum width=#1,
                     minimum height=0.85cm,
                     path picture={\draw[line width=0.5pt]
                       ([xshift=-1.1mm]path picture bounding box.south) --
                       ([yshift=1.6mm]path picture bounding box.south) --
                       ([xshift=1.1mm]path picture bounding box.south);}},
  registro/.default = 0.8cm,
  % condición de una cadena if / else-if: "sí" sale a la derecha, "no" abajo
  pregunta/.style = {draw=black, line width=0.5pt, rounded corners=7pt,
                     fill=white, font=\scriptsize\ttfamily, align=center,
                     inner sep=2.5pt, minimum height=0.62cm, text width=#1},
  pregunta/.default = 1.9cm,
  si/.style      = {font=\tiny, inner sep=1pt, above, pos=0.12},
  no/.style      = {font=\tiny, inner sep=1pt, left, pos=0.45},
  %
  titulo/.style  = {anchor=north west, font=\small, inner sep=4pt, align=left},
  rotulo/.style  = {anchor=north west, font=\scriptsize\itshape, black!70,
                    inner sep=2pt, align=left},
}

% submódulo: nombre en monoespaciada; lógica propia: nombre en romana
\newcommand{\blk}[2]{{\ttfamily\bfseries #1}\\[1pt]#2}
\newcommand{\etq}[2]{{\bfseries #1}\\[1pt]#2}
\newcommand{\punto}[1]{\fill (#1) circle (1.3pt);}

% Parámetros en el formato del datasheet: NOMBRE=valor en monoespaciada, una
% entrada por celda, hasta cuatro columnas. Una columna: \prm{A=1\\B=2};
% dos: \prm{A=1 & B=2\\C=3}.
\newcommand{\prm}[1]{{\ttfamily\fontsize{5.6}{6.6}\selectfont
  \begin{tabular}{@{}l@{\hspace{7pt}}l@{\hspace{7pt}}l@{\hspace{7pt}}l@{}}#1\end{tabular}}}

\def\anchopos{0.5}

% Mux de una decisión: trapecio con el lado ancho (las entradas) a la izquierda.
% #1 = nombre, #2 = centro, #3 = alto en cm. Deja tres coordenadas:
%   (#1-in)  centro del lado de entradas: las flechas llegan a (#1-in |- opcion.east)
%   (#1-out) la salida
%   (#1-sel) mitad del lado de abajo: por ahí entra el registro que elige
% El rótulo de adentro es opcional: \muxdec{m}{(4,0)}{3.6} y después
% \node[...] at (m) {...}.
\newcommand{\muxdec}[3]{%
  \path #2 coordinate (#1);
  \draw[line width=0.6pt, fill=grisr]
    ($(#1)+(-0.45,{#3/2})$) -- ($(#1)+(0.45,{#3/2-0.35})$) --
    ($(#1)+(0.45,{-#3/2+0.35})$) -- ($(#1)+(-0.45,{-#3/2})$) -- cycle;
  \coordinate (#1-in)  at ($(#1)+(-0.45,0)$);
  \coordinate (#1-out) at ($(#1)+(0.45,0)$);
  \coordinate (#1-sel) at ($(#1)+(0,{-#3/2+0.175})$);}

% El mismo mux rotado 90 grados, para un flujo de arriba hacia abajo: entradas
% por el lado ancho de arriba, salida abajo y el registro que elige por la
% izquierda. #1 = nombre, #2 = centro, #3 = ancho en cm. Deja (#1-in) centro del
% lado de entradas, (#1-out) la salida y (#1-sel) mitad del lado izquierdo.
\newcommand{\muxdecv}[3]{%
  \path #2 coordinate (#1);
  \draw[line width=0.6pt, fill=grisr]
    ($(#1)+({-#3/2},0.45)$) -- ($(#1)+({#3/2},0.45)$) --
    ($(#1)+({#3/2-0.35},-0.45)$) -- ($(#1)+({-#3/2+0.35},-0.45)$) -- cycle;
  \coordinate (#1-in)  at ($(#1)+(0,0.45)$);
  \coordinate (#1-out) at ($(#1)+(0,-0.45)$);
  \coordinate (#1-sel) at ($(#1)+({-#3/2+0.175},0)$);}

% Leyenda de colores + nota. #1 = posición, #2 = color/texto, #3 = nota.
\newcommand{\leyenda}[3]{%
  \begin{scope}[shift={(#1)}, every node/.style={font=\scriptsize}]
    \coordinate (t0) at (-0.5,0);
    \foreach \c/\t [count=\i] in {#2} {
      \pgfmathtruncatemacro\j{\i-1}
      \node[draw, fill=\c, minimum width=4mm, minimum height=3mm, inner sep=0pt,
            right=5mm of t\j.east, anchor=west] (l\i) {};
      \node[right=1.5mm of l\i] (t\i) {\t};
    }
    \node[anchor=north west, black!70, text width=19cm, align=left]
      at (-0.05,-0.3) {#3};
  \end{scope}}

% Lo que dicen los parámetros, para la nota de la leyenda
\newcommand{\notaparam}{\texttt{NOMBRE=valor}: parámetros con que el padre
  instancia el bloque (se omiten los que sólo pasan el del padre sin cambio).}

\tikzset{
  pre/.style     = {bloque=azul!35, param={azul!35}{azul!12}},
  varios/.style  = {double copy shadow={shadow xshift=2pt, shadow yshift=-2pt}},
  op/.style      = {oper, font=\scriptsize},
  caja/.style    = {desp, inner sep=2pt, minimum height=0.5cm},
  reg/.style     = {registro=#1, minimum height=0.55cm, font=\scriptsize\ttfamily},
  reg/.default   = 1.15cm,
  cnt/.style     = {registro=1.85cm, minimum height=0.5cm, font=\scriptsize\ttfamily},
  fondo f/.style = {draw=black!45, dashed, line width=0.5pt, rounded corners=3pt, fill=#1},
  rotf/.style    = {font=\scriptsize\itshape, black!75, inner sep=2pt},
  red/.style     = {font=\scriptsize\ttfamily, inner sep=1pt},
  chico/.style   = {font=\tiny, inner sep=1pt, align=center},
  y3/.style      = {logic gate input sep=0.2cm},
}
% \red[largo]{pin}{nombre}: la red que entra a un pin, por nombre, desde la izquierda
\newcommand{\red}[3][0.3]{\draw[flujo] ($(#2)+(-#1,0)$) node[left, red] {#3} -- (#2);}
% \redab[largo]{pin}{nombre}: la misma, desde abajo
\newcommand{\redab}[3][0.4]{\draw[flujo] ($(#2)+(0,-#1)$) node[below, red] {#3} -- (#2);}
% \redar[largo]{pin}{nombre}: la misma, desde arriba
\newcommand{\redar}[3][0.4]{\draw[flujo] ($(#2)+(0,#1)$) node[above, red] {#3} -- (#2);}
% \redsal[largo]{pin}{nombre}: la red que sale de un pin hacia la derecha
\newcommand{\redsal}[3][0.3]{\draw[flujo] (#2) -- ($(#2)+(#1,0)$) node[right, red] {#3};}

\begin{document}
\begin{tikzpicture}[fondo]

\def\xin{-0.3}    % donde arrancan las entradas
\def\xout{19.75}  % donde terminan las salidas
\def\ytop{14.5}   % de donde bajan las de configuración

% ============================================================ las regiones ===
\node[fondo f=amarillo!22, fit={(0.35,8.55) (5.75,13.55)}] (fdsp) {};
\node[fondo f=amarillo!22, fit={(6.15,8.55) (13.95,13.55)}] (fasg) {};
\node[fondo f=amarillo!22, fit={(14.35,8.55) (19.05,13.55)}] (fmar) {};
\node[fondo f=azul!10, fit={(0.35,5.15) (5.15,7.95)}] (fmue) {};
\node[fondo f=azul!10, fit={(0.35,-0.2) (5.15,4.65)}] (ftrp) {};
\node[fondo f=verde!14, fit={(5.55,-2.45) (7.95,0.7)}] (fumb) {};
\node[fondo f=grisr!45, fit={(11.85,1.75) (19.05,7.95)}] (fcon) {};
\node[fondo f=grisr!45, fit={(11.85,-0.75) (19.05,1.35)}] (fbas) {};
\node[rotf, anchor=north west] at (fdsp.north west) {disparo};
\node[rotf, anchor=north west] at (fasg.north west) {asignador};
\node[rotf, anchor=north west] at (fmar.north west) {marcas que se cargan al abrir};
\node[rotf, anchor=north west] at (fmue.north west) {la muestra, relativa a la base};
\node[rotf, anchor=north west] at (ftrp.north west) {banco de trapecios y región de la meseta};
\node[rotf, anchor=north west] at (fumb.north west) {ventana de amplitud};
\node[rotf, anchor=north west] at (fcon.north west) {contabilidad};
\node[rotf, anchor=north west] at (fbas.north west) {hacia la línea de base};

% =============================================================== disparo ===
% el flanco: trig_i en alto y trig_q todavía en bajo
\node[puerta y=ni, anchor=input 1] (edg) at (2.4,12.75) {};
\node[reg=0.95cm] (tq) at (1.45,11.8) {trig\_q};
\draw[flujo] (\xin,12.75) node[left, sen] {trig\_i} -- (edg.input 1);
\draw[flujo] (0.75,12.75) |- (tq.west);
\punto{0.75,12.75}
\draw[flujo] (tq.east) -| ($(edg.input 2)+(-0.2,0)$) -- (edg.input 2);
% flanco o pendiente
\node[puerta o, anchor=input 1] (oeo) at (4.0,0 |- edg.output) {};
\draw[flujo] (edg.output) -- (oeo.input 1);
\coordinate (je) at ($(edg.output)!0.35!(oeo.input 1)$);
\punto{je}
% pendiente: un flanco entre dos muestras válidas espera a la siguiente
\node[reg=1.25cm] (pend) at (2.85,9.65) {trig\_pend};
\node[chico, anchor=north] at (pend.south) {1: flanco $\wedge\neg$val;\; 0: val};
\draw[flujo] (je) -- (je |- 0,10.3) -| (pend.north);
\draw[flujo] (pend.east) -| ($(oeo.input 2)+(-0.25,0)$) -- (oeo.input 2);
% disparo = (flanco o pendiente), con muestra válida y en marcha
\node[puerta y=nnn, y3, anchor=input 1] (dsp) at (4.9,11.45) {};
\draw[flujo] (oeo.output) -- ++(0.15,0) |- (dsp.input 1);
\red{dsp.input 2}{val}
\red{dsp.input 3}{run}
\draw[flujo] (dsp.output) -- ++(0.12,0) -- ++(0,-0.6) node[below, red] {disparo};

% ============================================================= asignador ===
% habilitadas: las cfg_n_acc primeras
\node[caja] (hab) at (7.0,12.95) {las \texttt{n\_acc}\\ primeras};
\draw[ctrl] (7.0,\ytop) node[above, sen, black!80] {cfg\_n\_acc\_i} -- (hab.north);
% ocupada = no libre, o reservada en el ciclo anterior
\node[puerta o=in, anchor=input 1] (ocu) at (6.95,11.0) {};
\red{ocu.input 1}{libre}
\red{ocu.input 2}{reserva\_q}
% libre habilitada = habilitada y no ocupada
\node[puerta y=ni, anchor=input 2] (lib) at ($(ocu.output)+(0.35,0)$) {};
\draw[flujo] (ocu.output) -- (lib.input 2);
\draw[flujo] (hab.east) -| ($(lib.input 1)+(-0.2,0)$) -- (lib.input 1);
\node[chico, anchor=south] at ($(hab.east)+(0.32,0)$) {habil};
% puede: alguna libre; sel: la de menor índice
\node[op, anchor=west] (pue) at ($(lib.output)+(0.45,0)$) {$\neq0$};
\draw[flujo] (lib.output) -- (pue.west);
\coordinate (jl) at ($(lib.output)!0.4!(pue.west)$);
\punto{jl}
\draw[flujo] (pue.east) -- ++(0.5,0);
\node[chico, anchor=south] at ($(pue.east)+(0.25,0.02)$) {puede};
\node[caja, anchor=west] (pri) at ($(jl |- 0,9.85)+(0.2,0)$) {menor\\ índice};
\draw[flujo] (jl) |- (pri.west);
\node[caja, anchor=west] (oh) at ($(pri.east)+(0.2,0)$) {$1 \ll$ sel};
\draw[flujo] (pri.east) -- (oh.west);
% abre = disparo, sin veto, con alguna libre; perdido = disparo, sin veto, sin libre
\node[puerta y=nin, y3, anchor=input 1] (abr) at (11.2,12.15) {};
\red{abr.input 1}{disparo}
\red{abr.input 2}{veto}
\red{abr.input 3}{puede}
\node[puerta y=nii, y3, anchor=input 1] (per) at (11.2,13.25) {};
\red{per.input 1}{disparo}
\red{per.input 2}{veto}
\red{per.input 3}{puede}
% la ranura que abre: un solo bit, el de sel, si abre
\node[puerta y, anchor=input 2] (aoh) at ($(oh.east)+(0.3,0)$) {};
\draw[flujo] (oh.east) -- (aoh.input 2);
\draw[flujo] (abr.output) -- ++(0.2,0) coordinate (ja) -- (ja |- 0,10.55) -| ($(aoh.input 1)+(-0.15,0)$) -- (aoh.input 1);
\node[chico, anchor=west] at ($(ja)+(0.05,0.1)$) {abre};
% la decisión, registrada junto con la muestra
\node[reg=1.05cm] (rper) at (13.35,0 |- per.output) {perdido\_q};
\node[reg=1.05cm] (rabr) at (13.35,10.05) {abre\_q};
\node[reg=1.05cm] (rres) at (13.35,9.3) {reserva\_q};
\draw[flujo] (per.output) -- (rper.west);
\draw[flujo] (aoh.output) -- ++(0.12,0) coordinate (jo) |- (rabr.west);
\draw[flujo] (jo) |- (rres.west);
\punto{jo}
% esperan: alguna cerrada entre las habilitadas
\node[puerta y, anchor=input 1] (esp) at (7.3,9.05) {};
\red{esp.input 1}{cerrada}
\red{esp.input 2}{habil}
\node[op] (espn) at (8.6,0 |- esp.output) {$\neq0$};
\draw[flujo] (esp.output) -- (espn.west);
\draw[flujo] (espn.east) -- ++(0.3,0) node[right, red] {esperan};

% ================================================================ marcas ===
% muestras desde el disparo anterior: +1 por muestra válida, cero en el disparo
\node[reg=1.4cm] (des) at (15.35,12.55) {desde\_disp};
\node[op] (dmas) at (16.8,12.55) {$+1$};
\draw[flujo] (des.east) -- (dmas.west);
\draw[flujo] (dmas.north) |- ($(des.north)+(0,0.3)$) -- (des.north);
\node[chico, anchor=south west] at ($(des.north)+(0.05,0.32)$) {satura; 0 con el disparo};
\node[reg=1.3cm] (dtp) at (18.3,12.55) {dt\_prev\_q};
\draw[flujo] (dmas.east) -- (dtp.west);
\node[chico, anchor=north] at (dtp.south) {se carga con el disparo};
% cuántas había abiertas cuando se decidió
\node[op] (cue) at (16.8,11.0) {$\Sigma$};
\draw[flujo] ($(cue.west)+(-0.3,0)$) node[left, red, align=right] {abierta $\vee$\\ reserva\_q} -- (cue.west);
\node[reg=1.3cm] (otr) at (18.3,11.0) {otras\_q};
\draw[flujo] (cue.east) -- (otr.west);
\node[chico, anchor=north] at (otr.south) {satura a 4 b};
% el número de ventana, con cada apertura
\node[reg=1.4cm] (wid) at (15.35,9.4) {win\_id};
\node[op] (wmas) at (16.8,9.4) {$+1$};
\draw[flujo] (wid.east) -- (wmas.west);
\draw[flujo] (wmas.north) |- ($(wid.north)+(0,0.3)$) -- (wid.north);
\node[chico, anchor=south] at ($(wid.north)+(0,0.32)$) {con \texttt{abre}};
\node[reg=1.3cm] (wiq) at (18.3,9.4) {win\_id\_q};
\draw[flujo] (wmas.east) -- (wiq.west);

% ================================================================ muestra ===
\node[op] (res) at (1.85,6.65) {$-$};
\node[op] (max) at (3.0,6.65) {max};
\node[cte] (cer) at (3.0,5.9) {0};
\draw[flujo] (cer) -- (max.south);
\draw[flujo, ancho=14, ancho en=0.35] (\xin,6.65) node[left, sen] {dat\_i} -- (res.west);
\draw[flujo, ancho=14, ancho en=0.4] (\xin,5.6) node[left, sen] {baseline\_i} -- (res.south |- 0,5.6) -- (res.south);
\draw[flujo] (res.east) -- node[chico, above] {$x$} (max.west);
\node[reg=0.95cm, minimum height=1.15cm] (xq) at (4.4,6.35) {xc\_q\\ x\_q\\ val\_q};
\draw[flujo] (max.east) -- node[chico, above] {$x_c$} (xq.west |- max.east);
\coordinate (jx) at ($(res.east)!0.4!(max.west)$);
\draw[flujo] (jx) |- (xq.west |- 0,6.3);
\punto{jx}
\draw[flujo] (\xin,7.45) node[left, sen] {val\_i} -- (4.4,7.45) -- (xq.north);

% ===================================================== trapecios y meseta ===
\node[pre, text width=2.0cm, varios] (trp) at (2.75,3.6)
  {\blk{mca\_trap {\normalfont$\times$} 3}{}\nodepart{two}$k$, $l$, \texttt{sh} propios;\\ $M$ común};
\draw[flujo, anchov=15, ancho en=0.75] (4.4,0 |- xq.south) -- node[sen, right, pos=0.45] {x\_q} (4.4,3.9) -- (trp.east |- 0,3.9);
\draw[flujo] (\xin,3.6) node[left, sen, align=right] {cfg\_trap\_k/l\_i\\ cfg\_trap\_sh/mq\_i} -- (trp.west);
% la región de la meseta en el marco de la ranura: ini = off + 8, fin = ini + 2^lw
\node[op] (ini) at (2.0,1.95) {$+$};
\node[cte] (ocho) at (2.0,2.65) {8};
\draw[flujo] (ocho) -- (ini.north);
\draw[flujo] (\xin,1.95) node[left, sen] {cfg\_trap\_off\_i} -- (ini.west);
\node[caja] (pot) at (1.45,0.45) {$1 \ll \mathtt{lw}$};
\draw[flujo] (\xin,0.45) node[left, sen] {cfg\_trap\_lw\_i} -- (pot.west);
\coordinate (ji) at (2.75,1.95);
\node[op] (fin) at (2.75,1.2) {$+$};
\draw[flujo] (ji) -- (fin.north);
\draw[flujo] (pot.east) -| (fin.south);
\node[reg=0.85cm] (rini) at (4.4,1.95) {ini};
\node[reg=0.85cm] (rfin) at (4.4,1.2) {fin};
\draw[flujo] (ini.east) -- (rini.west);
\punto{ji}
\draw[flujo] (fin.east) -- (rfin.west);
\node[chico, anchor=north] at ($(rfin.south)+(0,-0.05)$) {uno por forma;\\ satura};

% ================================================== ventana de amplitud ===
\node[caja] (smin) at (6.2,0.05) {$\ll$ s};
\node[op] (mas1) at (6.2,-0.7) {$+1$};
\node[caja] (smax) at (6.2,-1.38) {$\ll$ s};
\node[op] (men1) at (6.2,-2.05) {$-1$};
\node[reg=0.85cm] (rmin) at (7.35,0.05) {min\_q};
\node[reg=0.85cm] (rmax) at (7.35,-2.05) {max\_q};
\draw[flujo] (smin.west) ++(-0.3,0) node[left, red] {min} -- (smin.west);
\draw[flujo] (mas1.west) ++(-0.3,0) node[left, red] {max} -- (mas1.west);
\draw[flujo] (mas1.south) -- (smax.north);
\draw[flujo] (smax.south) -- (men1.north);
\draw[flujo] (smin.east) -- (rmin.west);
\draw[flujo] (men1.east) -- (rmax.west);
\node[chico, anchor=north] at ($(rmin.south)+(0,-0.12)$) {una vez,\\ para las 3};

% ================================================================ ranuras ===
\node[bloque=verde!20, varios, text width=2.5cm, minimum height=9.7cm,
      path picture={\fill[verde!55] (path picture bounding box.north west) rectangle
        ([yshift=-0.5cm]path picture bounding box.north east);
        \draw[line width=0.6pt] ([yshift=-0.5cm]path picture bounding box.north west) --
        ([yshift=-0.5cm]path picture bounding box.north east);}]
  (slt) at (9.85,2.9) {};
\node[anchor=north, font=\scriptsize, inner sep=0pt] at ([yshift=-0.14cm]slt.north)
  {\texttt{\bfseries mca\_slot} $\times$ 3};
\node[font=\scriptsize, align=center, text width=2.5cm] at ([yshift=0.3cm]slt.center)
  {abre con \texttt{abre\_q}[$k$]\\ y guarda las marcas\\[8pt]
   cierra sola, cuando\\ \texttt{len} llega a la\\ \texttt{gate\_long} que\\ retuvo al abrir\\[8pt]
   queda cerrada hasta\\ \texttt{libera\_i}[$k$]};
% entradas: la muestra, los trapecios, sus ventanas y la de amplitud
\draw[flujo, ancho=14, ancho en=0.75] (xq.east |- 0,6.35) -- (slt.west |- 0,6.35);
\node[chico, anchor=south] at (6.2,6.4) {xc\_q, val\_q};
\draw[flujo] (trp.east |- 0,3.35) -- node[chico, above, pos=0.15] {3$\times$18} (slt.west |- 0,3.35);
\draw[flujo] (rini.east) -- (slt.west |- rini.east);
\draw[flujo] (rfin.east) -- (slt.west |- rfin.east);
\draw[flujo] (rmin.east) -- (slt.west |- rmin.east);
\draw[flujo] (rmax.east) -- (slt.west |- rmax.east);
% de arriba: la configuración y la decisión del asignador
\draw[ctrl, cruce] (\xin,4.9) node[left, sen, black!80, align=right] {cfg\_gate\_*\_i\\ cfg\_amp\_*\_i} -- (slt.west |- 0,4.9);
\draw[line width=0.8pt] (rper.east) -- (14.12,0 |- rper.east) -- (14.12,8.15);
\draw[line width=0.8pt] (rabr.east) -- (14.12,0 |- rabr.east);
\punto{14.12,0 |- rabr.east}
\draw[flujo] (14.12,8.15) -- (10.4,8.15) -- (10.4,0 |- slt.north);
\node[chico, anchor=south west] at (10.45,8.18) {abre\_q, perdido\_q, marcas};
\draw[line width=0.8pt] (18.3,0 |- wiq.south) -- (18.3,8.3) -- (14.12,8.3);
\punto{14.12,8.3}
% salidas: el resultado de cada ranura y su estado
\draw[flujo, ancho=3$\times$295, ancho en=0.08] (slt.east |- 0,-1.15) -- (\xout,-1.15) node[right, sen, align=left]
  {resultados\\[-1pt]\normalfont\tiny pico, $Q$, mesetas, marcas};
\draw[flujo] (slt.east |- 0,-1.75) -- (\xout,-1.75) node[right, sen] {acepta\_o};
\draw[ctrl] (\xout,-2.3) node[right, sen] {libera\_i} -- (slt.east |- 0,-2.3);
\draw[flujo] ($(slt.east)+(0,2.15)$) -- ++(0.3,0) node[right, red, align=left] {estado de\\ las ranuras};

% ============================================================ contabilidad ===
% los casilleros del tiempo muerto: sin ranura libre, y por qué
\node[puerta y=ii, anchor=input 1] (dwi) at (13.05,7.35) {};
\red{dwi.input 1}{puede}
\red{dwi.input 2}{esperan}
\node[puerta y=in, anchor=input 1] (ddi) at (13.05,6.45) {};
\red{ddi.input 1}{puede}
\red{ddi.input 2}{esperan}
\draw[flujo] (dwi.output) -- (\xout,0 |- dwi.output) node[right, sen] {dt\_win\_o};
\draw[flujo] (ddi.output) -- (\xout,0 |- ddi.output) node[right, sen] {dt\_div\_o};
\node[op] (nocu) at (13.3,5.55) {$\Sigma$};
\red{nocu.west}{ocupada}
\draw[flujo] (nocu.east) -- (\xout,0 |- nocu.east) node[right, sen] {n\_ocup\_o};
% cada disparo, a exactamente un contador
\node[cnt] (ccr) at (14.75,4.55) {CNT\_CROSS};
\node[cnt] (cop) at (14.75,3.95) {CNT\_OPEN};
\node[cnt] (clv) at (14.75,3.35) {LOST\_VETO};
\node[cnt] (cla) at (14.75,2.75) {LOST\_ACTIVE};
\node[cnt] (clb) at (14.75,2.15) {LOST\_BUSY};
\red{ccr.west}{disparo}
\red{cop.west}{abre}
\red{clv.west}{disparo$\wedge$veto}
\red{cla.west}{perdido$\wedge$dt\_win}
\red{clb.west}{perdido$\wedge$dt\_div}
% cada cierre: se suman, porque pueden cerrar varias en el mismo ciclo
\node[cnt] (cto) at (17.65,4.55) {CNT\_TOTAL};
\node[cnt] (crj) at (17.65,3.95) {REJ\_AMP};
\node[cnt] (cpi) at (17.65,3.35) {PILEUP};
\node[chico, anchor=north, align=left] at (17.65,3.0)
  {$+\Sigma$ cierra\\ $+\Sigma$ cierra $\wedge\,\neg$acepta\\ $+\Sigma$ cierra $\wedge$ solapada};
\node[chico, anchor=north] at (14.75,1.85) {cada uno: $+1$ por disparo};
\draw[ctrl, cruce] (16.2,7.7) node[above, sen, black!80] {cnt\_clr\_i} -- (16.2,0 |- ccr.north);
\draw[flujo] (cto.east) -- (\xout,0 |- cto.east) node[right, sen] {cnt\_*\_o};

% ============================================================ hacia la base ===
% se sigue la base sólo si ninguna mide, ninguna está reservada y nadie abre
\node[op] (nab) at (13.1,0.75) {$\neq0$};
\red{nab.west}{abierta}
\node[op] (nre) at (13.1,0.05) {$\neq0$};
\red{nre.west}{reserva\_q}
\node[puerta y=iii, y3, anchor=input 1] (trk) at (14.1,0.75) {};
\draw[flujo] (nab.east) -- (trk.input 1);
\draw[flujo] (nre.east) -| ($(trk.input 2)+(-0.2,0)$) -- (trk.input 2);
\draw[flujo] ($(trk.input 3)+(-0.42,-0.42)$) node[below, red] {abre} |- (trk.input 3);
\draw[flujo] (trk.output) -- (\xout,0 |- trk.output) node[right, sen] {bl\_track\_o};
\node[op] (ncie) at (16.55,0.05) {$\neq0$};
\red{ncie.west}{cierra}
\draw[flujo] (ncie.east) -- (\xout,0 |- ncie.east) node[right, sen] {cierra\_o};
\node[op] (nest) at (16.55,-0.5) {$\neq0$};
\red{nest.west}{estancada}
\draw[flujo] (nest.east) -- (\xout,0 |- nest.east) node[right, sen] {bl\_stale\_set\_o};

% ============================================================== el límite ===
\node[contenedor, fit={(0.15,13.8) (19.3,-2.75)}] (lim) {};
\node[titulo, anchor=south west, inner sep=2pt, align=left] at (lim.north west)
  {\texttt{\bfseries mca\_window}\;\; {\scriptsize\itshape expandido}\\
   \prm{N\_ACC=3 & N\_FORMAS=3 & DW=14 & QW=32\\ LEN\_W=16 & AMP\_W=16 & PUB\_IDW=16 & NW=4\\
        TW=18 & TNW=8 & TDAW=10 & TMW=20}};
\node[sen, anchor=east, align=right] at (\xin,13.55)
  {veto\_i $\to$ veto\\ cfg\_run\_i $\to$ run};

% ================================================================= leyenda ===
\leyenda{0.2,-3.5}{amarillo!22/segmentación, azul!10/línea de base y conformado,
  verde!14/estadísticos, grisr!45/contabilidad, grisr/registro}{}

\end{tikzpicture}
\end{document}
